\section{Условие}
{\bfseries Цель работы:} изучение системных вызовов операционных систем Windows и Linux для создания
процессов и анонимных каналов, а также приобретение навыков построения кроссплатформенной
обёртки (единого API) над системными примитивами.

{\bfseries Задание:} разработать программу, состоящую из двух процессов:
\begin{itemize}
  \item родительский процесс \texttt{parent} считывает со стандартного ввода строки,
    содержащие числа с плавающей запятой, и передаёт их дочернему процессу по первому
    анонимному каналу;
  \item дочерний процесс \texttt{child} получает числа со стандартного ввода, накапливает
    их сумму, записывает отчёт в указанный в аргументе файл и возвращает итоговую сумму
    родителю по второму анонимному каналу;
  \item родительский процесс выводит полученную сумму и ожидает завершения дочернего процесса.
\end{itemize}
Программа должна собираться и выполняться одинаково в средах Windows и Linux: логика приложений
не должна содержать системных вызовов, всё платформозависимое должно быть скрыто за единым
кроссплатформенным API, выбор реализации (Win32 или POSIX) должен выполняться на этапе сборки
проекта (CMake).

{\bfseries Вариант:} 2